%! AP Statistics — Unit 2, Lesson 6
\documentclass[10pt]{article}
\usepackage[margin=0.6in]{geometry}
\usepackage{xcolor}
\usepackage[most]{tcolorbox}
\usepackage{enumitem}
\usepackage{multicol}
\usepackage{amsmath,amssymb}
\usepackage{array}
\usepackage{tabularx}
\tcbuselibrary{breakable}

\definecolor{sky}{HTML}{EAF3FB}
\definecolor{navy}{HTML}{1F3A5F}
\definecolor{redbox}{HTML}{FCECEC}
\definecolor{greenbox}{HTML}{EDF8F0}
\definecolor{goldbox}{HTML}{FFF7E6}
\newtcolorbox{skillbox}[2][]{
    colback=#2, colframe=navy, boxrule=0.8pt, arc=2mm,
    left=2mm,right=2mm,top=1.5mm,bottom=1.5mm,
    fonttitle=\bfseries, title={#1}, halign title=center, breakable
}
\setlist[itemize]{leftmargin=1.2em,itemsep=0.35em,topsep=0.35em}
\setlist[enumerate]{leftmargin=1.4em,itemsep=0.35em,topsep=0.35em}
\newcolumntype{C}[1]{>{\centering\arraybackslash}p{#1}}
\newcolumntype{Y}{>{\centering\arraybackslash}X}

\newcommand{\CourseName}{AP Statistics}
\newcommand{\SchoolYear}{2026--2027}
\newcommand{\MeetingLength}{55 minutes}
\newcommand{\UnitNumberName}{Unit 2: Exploring Two-Variable Data \quad}
\newcommand{\LessonNumberName}{Lesson 2.6: Linear Regression Models}

\begin{document}
\begin{center}
    {\Large\bfseries \CourseName: \SchoolYear\ (\MeetingLength) \\
    {\small\bfseries \UnitNumberName \LessonNumberName}}
\end{center}

\begin{tcolorbox}[colback=sky,colframe=navy,boxrule=0.9pt,arc=2mm,
    left=2mm,right=2mm,top=1.5mm,bottom=1.5mm]
    \textbf{Primary Objective:} Students will write and interpret a linear regression equation
    $\hat{y} = a + bx$, including interpretation of the slope and $y$-intercept in context.
    Students will use the equation to make predictions and will understand the limitations of
    predictions outside the data range (Big Idea: DAT; AP Skill 2).
\end{tcolorbox}

\begin{skillbox}[Priority Ideas \& Skills]{goldbox}
\begin{minipage}[t]{0.48\linewidth}
    \textbf{AP Skill 2 --- Data Analysis}
    \begin{itemize}
        \item Write the regression equation in the form $\hat{y} = a + bx$ with variable names (not generic $x$ and $y$).
        \item Interpret the slope: ``For each additional [one unit of $x$], the predicted [y] [increases/decreases] by [slope value].''
        \item Interpret the $y$-intercept in context (and recognize when it is not meaningful).
        \item Use the equation to make predictions; identify interpolation vs.\ extrapolation.
    \end{itemize}
\end{minipage}
\hfill
\begin{minipage}[t]{0.48\linewidth}
    \textbf{Key Understandings}
    \begin{itemize}
        \item $\hat{y}$ (``$y$-hat'') is a \textit{predicted} value, not an exact one; actual values differ from predictions.
        \item The slope is the predicted change in $y$ per one-unit change in $x$ --- always include units and direction.
        \item Extrapolation (predicting outside the data range) is unreliable and should be flagged.
        \item A regression line always passes through $(\bar{x}, \bar{y})$.
    \end{itemize}
\end{minipage}
\end{skillbox}

\begin{skillbox}[{Vocabulary, Concepts, \& Theorems}]{greenbox}
\begin{minipage}[t]{0.48\linewidth}
    \begin{itemize}
        \item \textbf{Regression line (least-squares line)}: $\hat{y} = a + bx$; the line that best summarizes the linear relationship
        \item \textbf{Slope} $b$: predicted change in $\hat{y}$ for a one-unit increase in $x$
        \item \textbf{$y$-intercept} $a$: predicted value of $y$ when $x = 0$; not always meaningful in context
        \item \textbf{$\hat{y}$ (predicted value)}: the value of $y$ estimated from the regression equation for a given $x$
    \end{itemize}
\end{minipage}
\hfill
\begin{minipage}[t]{0.48\linewidth}
    \begin{itemize}
        \item \textbf{Interpolation}: predicting $\hat{y}$ for an $x$ value within the range of the data; generally reliable
        \item \textbf{Extrapolation}: predicting $\hat{y}$ for an $x$ value outside the range of the data; unreliable; should be avoided or flagged
        \item \textbf{Residual}: actual $y$ minus predicted $\hat{y}$; $e = y - \hat{y}$ (detail in Lesson 2.7)
        \item \textbf{Point $(\bar{x}, \bar{y})$}: the regression line always passes through the mean of $x$ and mean of $y$
    \end{itemize}
\end{minipage}
\end{skillbox}

\begin{skillbox}[Activate Prior Knowledge \& Spiral Review]{sky}
\begin{minipage}[t]{0.48\linewidth}
    \textbf{Quick Review (3 min)}\\
    Students recall slope-intercept form $y = mx + b$ from Algebra. Ask: ``What does the slope mean in a real-world context?'' Bridge: in statistics, the slope answers ``how does $y$ change, on average, for each one-unit increase in $x$?'' The hat on $\hat{y}$ signals it is a prediction.
\end{minipage}
\hfill
\begin{minipage}[t]{0.48\linewidth}
    \textbf{Spiral Connections}
    \begin{itemize}
        \item Lesson 2.5: $r$ described the strength of the linear association; regression uses that association to make predictions.
        \item Lesson 2.7: the difference between $y$ and $\hat{y}$ is the residual.
        \item Lesson 2.8: the formulas for $a$ and $b$ come from minimizing the sum of squared residuals.
    \end{itemize}
\end{minipage}
\end{skillbox}

\begin{skillbox}[Hook \& Lesson]{sky}
\begin{minipage}[t]{0.48\linewidth}
    \textbf{Hook (5 min)}\\
    Display the scatterplot of height vs.\ shoe size. Ask: ``If I told you a person's height was 68 inches, what would you predict for their shoe size?'' Students informally draw a trend line and read off a prediction. Today we formalize that line and make precise predictions with an equation.
\end{minipage}
\hfill
\begin{minipage}[t]{0.48\linewidth}
    \textbf{Lesson Flow (15 min)}
    \begin{enumerate}
        \item Introduce the regression equation $\hat{y} = a + bx$ with specific variable names.
        \item Walk through slope and $y$-intercept interpretation for a concrete example.
        \item Make a prediction: substitute an $x$ value, compute $\hat{y}$, interpret.
        \item Identify an extrapolation: flag it, explain the risk.
        \item Verify: does the line pass through $(\bar{x}, \bar{y})$?
    \end{enumerate}
\end{minipage}
\end{skillbox}

\begin{skillbox}[Explicit Instruction \& Active Monitoring]{sky}
\begin{minipage}[t]{0.48\linewidth}
    \textbf{Direct Instruction (10 min)}
    \begin{itemize}
        \item AP slope interpretation format: ``For each additional [unit of $x$], the predicted [y-variable] [increases/decreases] by approximately [|slope| with units].''
        \item $y$-intercept caveat: if $x = 0$ is outside the data range or physically impossible, state that the intercept ``is not meaningful in context.''
        \item Prediction: $\hat{y}$ is a \textit{predicted average} --- individual values will vary.
    \end{itemize}
\end{minipage}
\hfill
\begin{minipage}[t]{0.48\linewidth}
    \textbf{Active Monitoring}
    \begin{itemize}
        \item Common AP error: interpreting the slope without units or direction (``slope = 2.3'' with no context earns no credit).
        \item Probe: ``Is $x = 0$ in the data range? Can a person have 0 hours of sleep?'' --- tests whether the intercept is meaningful.
        \item Check: does the student use $\hat{y}$ (not $y$) when writing predictions?
    \end{itemize}
\end{minipage}
\end{skillbox}

\begin{skillbox}[Group Work \& Differentiation]{redbox}
\begin{minipage}[t]{0.48\linewidth}
    \textbf{Group Activity (8--10 min)}\\
    Groups are given a regression equation from technology output (e.g., $\widehat{\text{price}} = 24500 - 1800(\text{age})$ for used cars). Groups:
    \begin{enumerate}
        \item Interpret the slope in context.
        \item Interpret the $y$-intercept and assess its meaning.
        \item Predict price for a 3-year-old and a 20-year-old car. Flag any extrapolation.
        \item Identify one limitation of using this model.
    \end{enumerate}
\end{minipage}
\hfill
\begin{minipage}[t]{0.48\linewidth}
    \textbf{Differentiation}
    \begin{itemize}
        \item \textit{Support}: Provide fill-in interpretation templates: ``For each additional [year], the predicted price [decreases] by \$\underline{\hspace{1cm}}.''
        \item \textit{Extension}: Students generate two predictions for the same $x$ from two different regression equations (different datasets) and explain why the predictions differ.
        \item \textit{Technology}: Input data into a graphing calculator or Desmos and generate the regression equation; confirm it matches the given equation.
    \end{itemize}
\end{minipage}
\end{skillbox}

\begin{skillbox}[Individual Work \& Assessment]{redbox}
\begin{minipage}[t]{0.48\linewidth}
    \textbf{Exit Ticket (5 min)}\\
    Regression output: $\widehat{\text{score}} = 40 + 5(\text{hours studied})$. Students: (1) interpret slope, (2) interpret $y$-intercept, (3) predict score for 6 hours of studying, (4) explain why predicting for 30 hours is problematic.
\end{minipage}
\hfill
\begin{minipage}[t]{0.48\linewidth}
    \textbf{AP-Style Check}\\
    \textit{``The regression equation for predicting weight from height is $\widehat{\text{weight}} = -105 + 3.7(\text{height})$. Interpret the slope and $y$-intercept in context. Is the $y$-intercept meaningful? Explain.''}\\[4pt]
    Full AP credit requires: units, direction, context for slope; recognition that weight $= 0$ at height $= 0$ is not meaningful.
\end{minipage}
\end{skillbox}

\begin{skillbox}[Reinforcement \& Extension]{goldbox}
\begin{minipage}[t]{0.48\linewidth}
    \textbf{Homework / Practice}
    \begin{itemize}
        \item Given 3 regression equations in different contexts, interpret slope and intercept, make a prediction, and identify any extrapolation.
        \item AP textbook FRQ practice on regression interpretation.
    \end{itemize}
\end{minipage}
\hfill
\begin{minipage}[t]{0.48\linewidth}
    \textbf{Extension / Preview}
    \begin{itemize}
        \item \textit{Extension}: Compare two regression equations for the same response variable using different predictors. Which model would you prefer, and why?
        \item \textit{Preview}: Lesson 2.7 --- the regression line doesn't pass through every point. The difference between the actual and predicted value is called a \textit{residual}. Positive residuals mean we underestimated; negative residuals mean we overestimated.
    \end{itemize}
\end{minipage}
\end{skillbox}

\end{document}
